% Generated draft: requires mathematical and source review.
\documentclass[11pt]{article}
\usepackage[margin=25mm]{geometry}
\usepackage{fontspec}
\setmainfont{Latin Modern Roman}
\usepackage{amsmath,amssymb,mathtools,unicode-math}
\setmathfont{Latin Modern Math}
\usepackage{enumitem,needspace,xurl,hyperref}
\setlength{\parindent}{0pt}
\setlength{\parskip}{5pt}
\setlength{\emergencystretch}{3em}
\newcommand{\N}{\mathbb{N}}
\newcommand{\Z}{\mathbb{Z}}
\newcommand{\Q}{\mathbb{Q}}
\newcommand{\R}{\mathbb{R}}
\newcommand{\C}{\mathbb{C}}
\newcommand{\F}{\mathbb{F}}
\newcommand{\A}{\mathbb{A}}
\newcommand{\PP}{\mathbb{P}}
\newcommand{\E}{\mathbb{E}}
\newcommand{\eps}{\varepsilon}
\newcommand{\dd}{\,\mathrm{d}}
\begin{document}
\section*{UnsolvedMath ID: 30002619}
\textbf{Source problem code:} catalogue entry
\subsection{Definitions and mathematical statement}
Let \(\mathcal L^3\) denote Lebesgue measure on \(\R^3\). A \(2\)-current \(T\) is normal if both \(T\) and its boundary \(\partial T\) have finite mass, where \(\partial T(\eta)=T(\dd\eta)\) for every smooth compactly supported \(1\)-form \(\eta\). Write \(N_2(\R^3)\) for this class.

Retain exactly the vector fields supplied in the record:
\[
X(x,y,z)=(1,0,-x/2),\qquad
Y(x,y,z)=(0,1,y/2),\qquad
\xi=X\wedge Y.
\]
The source does not specify whether \(\mu\) is a nonnegative or signed scalar Radon measure. Under either interpretation, interpret its notation \(T=\xi\wedge\mu\) as the current represented by
\[
T(\omega)=\int_{\R^3}\langle\omega(p),\xi(p)\rangle\,\dd\mu(p)
\]
for smooth compactly supported \(2\)-forms \(\omega\). The hypothesis \(T\in N_2(\R^3)\) includes the required global finite-mass conditions. Absolute continuity \(\mu\ll\mathcal L^3\) means that \(\mu(A)=0\) for every Borel set \(A\) with \(\mathcal L^3(A)=0\).

The literal measure question in the supplied record is whether, for every such \(\mu\),
\[
\xi\wedge\mu\in N_2(\R^3)
\quad\Longrightarrow\quad
\mu\ll\mathcal L^3.
\]
\emph{Unresolved formulation conflict.} The record calls \(X,Y\) horizontal fields of the Heisenberg group \(\mathbb H^1\) and calls their span non-involutive. However, the displayed fields satisfy \([X,Y]=0\). Since they are everywhere linearly independent, their span is a smooth involutive rank-two distribution. Thus the literal displayed question and the described Heisenberg non-involutive target are incompatible. The implication above records the supplied formulas; it does not repair or identify the intended non-involutive example.

The record also asks for a classification of normal \(k\)-currents that fail a Frobenius-type involutivity conclusion asserted for integral currents, and asks what survives below \(C^1\) regularity or when involutivity is defined without Lie brackets. These are retained as research questions only: the supplied record does not state the precise integral-current conclusion, the proposed classification class, the ambient dimensions and hypotheses, the weaker regularity class, or the replacement definition of involutivity. Consequently, these extensions do not yet determine precise quantified mathematical statements.
\subsection{Short English statement}
The supplied question asks whether normality of a current represented by a prescribed plane field forces its representing measure to be absolutely continuous. Its displayed fields commute, conflicting with the intended non-involutive Heisenberg example. The broader classification and weaker-regularity questions lack necessary specifications. The supplied partially solved label is dated curation, not an independently verified status.
\subsection{Sources}
\begin{itemize}
\item Local UnsolvedMath registry: supplied complete source record for permanent ID \(30002619\), source problem code \texttt{catalogue entry}, title \emph{Involutivity Obstructions for Normal Currents}; registry paths \texttt{data/unsolvedmath/problems.json} and \texttt{catalog/problems.json}. These paths are registry references supplied by the project instructions; their contents were not independently inspected for this draft.
\item Supplied original report citation: \emph{Calculus of Variations (2015)}, Oberwolfach Reports, DOI \texttt{10.4171/owr/2014/33}, pp.\ 1801--1860. \url{https://doi.org/10.4171/owr/2014/33}. The record gives proposed year 2014 and no named proposer. The citation's year 2015 and the record's 2014 dating require reconciliation.
\item Supplied literature reference: Alberti--Massaccesi--Stepanov, ``currents tangent to distributions.'' \url{https://doi.org/10.4310/jdg/1668186786}. This is the descriptive attribution supplied in the record; no fuller bibliographic details were supplied.
\item Supplied literature reference: ``original Oberwolfach report.'' \url{https://doi.org/10.4171/owr/2014/33}.
\item Literature assessment dated exactly \texttt{2026-08-21}: the supplied curation labels the problem \texttt{partially\_solved} and states that later work proves substantial absolute-continuity and structure results for normal currents tangent to smooth distributions, while weaker-regularity and weaker-involutivity variants explicitly remain unresolved. This is supplied dated curation, not independently verified literature assessment; the cited claims and time-sensitive status require independent verification.
\item The record labels its statement \texttt{corrected\_verified} and says that notation, typography, or omitted definitions were repaired while preserving mathematical content. This is a supplied editorial qualification, not mathematical verification performed for this draft.
\end{itemize}
\end{document}
